\documentclass[conference]{IEEEtran}
\IEEEoverridecommandlockouts
\usepackage{cite}
\usepackage{amsmath,amssymb,amsfonts}
\usepackage{algorithmic}
\usepackage{graphicx}
\usepackage{textcomp}
\usepackage{xcolor}

\begin{document}

\title{NRP-1: A 0.47\,mW Event-Driven Neuromorphic Reflex Processor with On-Chip STDP Learning for Sub-2\,ms Edge Actuation}

\author{\IEEEauthorblockN{ECE Undergraduate Major Research Capstone Team}
\IEEEauthorblockA{\textit{Department of Electronics and Communication Engineering} \\
\textit{Undergraduate Major Project Thesis}\\
Email: capstone@ece.edu}
}

\maketitle

\begin{abstract}
Traditional edge computing architectures execute static neural network models on power-hungry microcontrollers or micro-GPUs, incurring 15--45\,ms software latencies and 30--85\,mW dynamic power while lacking local adaptivity to mechanical drift. This paper presents NRP-1, a custom event-driven Neuromorphic Reflex Processor implemented in synthesizable IEEE 1364-2001 Verilog and qualified for SkyWater 130nm CMOS. By combining 8-channel asynchronous delta-spike modulation, 256 plastic synapses with on-chip Spike-Timing-Dependent Plasticity (STDP), 32 integer Leaky Integrate-and-Fire (LIF) neurons, and an autonomous 4\,kHz SPI telemetry master, NRP-1 delivers a deterministic sub-2\,ms physical reflex ($1.601\,$ms verified) at $0.47\,$mW FPGA core power and $320\,\mu$W ASIC power. Evaluated on a variable-speed motor vibration test vehicle accelerating from 50\,Hz to 120\,Hz, NRP-1 achieves 0.00\% false alarms while fixed second-order high-pass filters trigger false trips as speed harmonics enter the passband. The architecture was formally verified via a 7-signal per-sample cycle-accurate lockstep engine across 10,000 steps with zero discrepancies between RTL and golden simulation.
\end{abstract}

\begin{IEEEkeywords}
Neuromorphic computing, Spiking Neural Networks, STDP, synthesizable Verilog, digital VLSI, cyber-physical reflexes.
\end{IEEEkeywords}

\section{Introduction}
Modern cyber-physical systems require sub-millisecond reflexes to prevent catastrophic mechanical failure...

\end{document}
